\section{Proofs}\label{app:proofs}

Throughout, $K=k_m$ is the largest grid point. Exact Pareto tails ($\delta\equiv0$) are not
in the class of Assumption~\ref{ass:A1}; wherever we argue ``under $p^{\otimes K}$'' the argument is
for exact Pareto excesses directly and does not use Assumption~\ref{ass:A1}. Further, $E\sim\mathrm{Exp}(1)$, and for $s<1$ we use
$\E e^{sE}=(1-s)^{-1}$ and $\E(Ee^{sE})=(1-s)^{-2}$.

\subsection{Scores, efficiency and the bias factor}

\begin{lemma}\label{lem:scores}
Let $Y$ have the EPD$(\xi,\delta;\rho')$ distribution and $E=\xi^{-1}\log Y$. At $\delta=0$:
\begin{enumerate}
\item[(i)] the scores for $(\xi,\delta)$ are $\xi^{-1}(E-1)$ and $\xi^{-1}a_{\rho'}(E)$, with
$a_{\rho'}(E)=(1-\rho')e^{\rho'E}-1$, and the Fisher information is
$\xi^{-2}\begin{pmatrix}1&c'\\c'&d'\end{pmatrix}$ with $c'=\rho'/(1-\rho')$,
$d'=\rho'^2/(1-2\rho')$; its determinant is $\xi^{-4}\rho'^4/\{(1-2\rho')(1-\rho')^2\}>0$;
\item[(ii)] the efficient influence function for $\xi$ is $\xi\psi_{\rho'}(E)$ with $\psi_{\rho'}$
as in Section~\ref{sec:bvm}, $\E\psi_{\rho'}(E)=0$ and $\operatorname{var}\psi_{\rho'}(E)=s(\rho')^2$;
\item[(iii)] differentiating $y^{-1/\xi}\{1+\xi^{-1}\delta(y^{\rho/\xi}-1)\}$, the conditional
survival function implied by Assumption~\ref{ass:A1} to first order in $\delta$, at
$\delta=0$ gives the log-density derivative $h_\rho$ of \eqref{eq:h}, and
$h_\rho(Y)=\xi^{-1}a_\rho(E)$;
\item[(iv)] $\operatorname{cov}\{\psi_{\rho'}(E),a_\rho(E)\}=\beta(\rho,\rho')$ and
$\operatorname{cov}\{E-1,a_\rho(E)\}=\rho/(1-\rho)$;
\item[(v)] under exact Pareto excesses, the estimator of \citet{bjs2009} at the working
value $\rho'$ satisfies $\sqrt{k}\{\hat\xi(\rho')-\xi\}/\xi=k^{-1/2}\sum_{i\le k}\psi_{\rho'}(E_i)+o_P(1)$.
\end{enumerate}
\end{lemma}

\begin{proof}
(i) The EPD log-density at $\delta=0$ is the Pareto log-density, whose $\xi$-score is
$\xi^{-1}(E-1)$; the dependence of $\rho'/\xi$ on $\xi$ does not enter at $\delta=0$.
Differentiating the EPD log-density in $\delta$ at $\delta=0$ gives $\xi^{-1}a_{\rho'}(E)$.
The moments $\operatorname{var}(E)=1$, $\operatorname{cov}(E,a_{\rho'})=\rho'/(1-\rho')$ and
$\operatorname{var}(a_{\rho'})=\rho'^2/(1-2\rho')$ follow from the two exponential moments.
(ii) The efficient influence function is the first row of the inverse information applied to
the scores: $\xi\{d'(E-1)-c'a_{\rho'}(E)\}/(d'-c'^2)$, and
$d'/(d'-c'^2)=(1-\rho')^2/\rho'^2$, $c'/(d'-c'^2)=(1-2\rho')(1-\rho')/\rho'^3$. Its variance is
$\xi^2d'/(d'-c'^2)=\xi^2s(\rho')^2$.
(iii) Direct differentiation; with $y=e^{\xi E}$ the result is $\xi^{-1}\{(1-\rho)e^{\rho E}-1\}$.
(iv) $\operatorname{cov}(a_{\rho'},a_\rho)=\rho\rho'/(1-\rho-\rho')$, and
$\{d'\rho/(1-\rho)-c'\rho\rho'/(1-\rho-\rho')\}/(d'-c'^2)=\beta(\rho,\rho')$ by the identity
$(1-\rho')(1-\rho-\rho')-(1-2\rho')(1-\rho)=\rho'(\rho'-\rho)$.
(v) The estimator is $\hat\xi=H-\hat\delta\rho'/(1-\rho')$ with $H=k^{-1}\sum\xi E_i$,
$\hat\delta=H(1-2\rho')(1-\rho')^3\rho'^{-4}\{M-(1-\rho')^{-1}\}$ and
$M=k^{-1}\sum\exp(\rho'\xi E_i/H)$. Under exact Pareto excesses the $E_i$ are independent
standard exponential and $\varepsilon=H/\xi-1=O_P(k^{-1/2})$. Write
$M=k^{-1}\sum f(E_i,\varepsilon)$ with $f(e,t)=\exp\{\rho'e/(1+t)\}$. As $\rho'<0$ and $e\ge0$,
$|\partial_t^2f(e,t)|\le C(e+e^2)$ for $|t|\le1/2$, so
$M=k^{-1}\sum e^{\rho'E_i}-\varepsilon\rho'k^{-1}\sum E_ie^{\rho'E_i}+O_P(k^{-1})$.
Up to $o_P(k^{-1/2})$, $\hat\xi$ is therefore a smooth function of three sample means with finite
variances, and the delta method gives an asymptotically linear expansion whose
influence function, after the algebra, equals $\psi_{\rho'}$. \qedhere
\end{proof}

\begin{proof}[Proof of Lemma~\ref{lem:beta}]
All factors of $\beta$ other than $\rho'-\rho$ are positive or pair to the positive ratio
$\rho/\rho'$. The Hill ratio follows from (iv): the Hill estimator has influence $E-1$, so its
bias is $\lambda\rho/(1-\rho)$ by the argument of Theorem~\ref{thm:bvm}(ii). \qedhere
\end{proof}

\begin{lemma}\label{lem:suff}
If $q_u/p=1+\delta(u)h_\rho+r_u$ with $\sup_{y\ge1}|r_u(y)|y^{-\epsilon}=o\{\delta(u)\}$ for some
$\epsilon<1/(4\xi)$, then Assumption~\ref{ass:A2} holds.
\end{lemma}

\begin{proof}
Let $v=\delta h_\rho+r_u\ge-1$ and $g(v)=(1+v)^{1/2}-1-v/2$, so that $|g(v)|\le v^2/2$ for
$v\ge-1$. The integrand of Assumption~\ref{ass:A2} is $p\{g(v)+r_u/2\}^2\le2pg(v)^2+pr_u^2/2$.
Since $h_\rho$ is bounded on $[1,\infty)$ and $\int y^{4\epsilon}p(y)\,dy<\infty$ for
$4\epsilon<1/\xi$, $\int pr_u^2=o(\delta^2)$ and
$\int pg(v)^2\le\frac14\int pv^4\le2\int p(\delta^4h_\rho^4+r_u^4)=o(\delta^2)$. \qedhere
\end{proof}

\subsection{Joint limits under Pareto tails}

\begin{lemma}\label{lem:clt}
Let $E_1,\dots,E_K$ be independent standard exponential and, for $j\le m$ and $i\le k_j$,
$V^{(j)}_i=\sum_{l=i}^{k_j}E_l/l$. For functions $\phi$ with $\E\phi(E)=0$, $\E\phi(E)^2<\infty$ and
$\phi',\phi''$ bounded on $[0,\infty)$, put $S_j(\phi)=k_j^{-1/2}\sum_{i\le k_j}\phi(V^{(j)}_i)$.
Then for finitely many such pairs $(j,\phi)$ the vector of the $S_j(\phi)$ is asymptotically
centred normal, and for two such functions $\phi,\chi$ and $t_j\le t_{j'}$,
$\lim\operatorname{cov}\{S_j(\phi),S_{j'}(\chi)\}=(t_j/t_{j'})^{1/2}\operatorname{cov}\{\phi(E),\chi(E+\log(t_{j'}/t_j))\}$.
\end{lemma}

Under exact Pareto excesses above $X_{n-K,n}$, R\'enyi's representation gives
\[
\xi^{-1}\log(X_{n-i+1,n}/X_{n-k_j,n})=V^{(j)}_i\qquad(i\le k_j\le K)
\]
jointly, with the same $E_l$, namely the normalised spacings $l\,\xi^{-1}\log(X_{n-l+1,n}/X_{n-l,n})$.
So every statistic $k_j^{-1/2}\sum_{i\le k_j}\phi\{\xi^{-1}\log(X_{n-i+1,n}/X_{n-k_j,n})\}$ equals
$S_j(\phi)$ in law, jointly in $j$ and $\phi$. The functions $\psi_{\rho'}$ and $a_\rho$ satisfy the conditions of the
lemma, because $\rho,\rho'<0$ make $e^{\rho z}$ and its derivatives bounded on $[0,\infty)$.

\begin{proof}
\emph{Linearisation.} Fix $j$ and write $r_i=\sum_{l=i}^{k_j}1/l=\E V^{(j)}_i$ and
$V^{(j)}_i-r_i=\sum_{l=i}^{k_j}(E_l-1)/l$, so that $\E(V^{(j)}_i-r_i)^2=\sum_{l\ge i}l^{-2}\le2/i$.
Taylor's theorem with $|\phi''|\le C$ gives
\[
S_j(\phi)=k_j^{-1/2}\sum_{i\le k_j}\phi(r_i)+k_j^{-1/2}\sum_{l\le k_j}w_{jl}(E_l-1)+R_j,\qquad
w_{jl}=\frac1l\sum_{i\le l}\phi'(r_i),
\]
with $|R_j|\le Ck_j^{-1/2}\sum_{i\le k_j}(V^{(j)}_i-r_i)^2$. Since
$\E(V^{(j)}_i-r_i)^4\le9(\sum_{l\ge i}l^{-2})^2\le36/i^2$, Minkowski's inequality gives
$(\E R_j^2)^{1/2}\le Ck_j^{-1/2}\sum_{i\le k_j}6/i=O(k_j^{-1/2}\log k_j)\to0$. The weights satisfy
$|w_{jl}|\le\sup|\phi'|$. The deterministic term is a Riemann sum: $r_i=\log(k_j/i)+O(1/i)$,
so $k_j^{-1}\sum_{i\le k_j}\phi(r_i)=\int_0^1\phi(-\log s)\,ds+O(k_j^{-1}\log k_j)$, using
$|\phi(r_i)-\phi(\log(k_j/i))|\le C/i$ and the monotone comparison of the sum with the integral
for the function $s\mapsto\phi(-\log s)$, whose variation on $[1/k_j,1]$ is $O(\log k_j)$, together with
$\int_0^{1/k_j}|\phi(-\log s)|\,ds=O(k_j^{-1}\log k_j)$, as $|\phi(z)|\le C(1+z)$. Since
$\int_0^1\phi(-\log s)\,ds=\E\phi(E)=0$, the deterministic term is $O(k_j^{-1/2}\log k_j)\to0$.

\emph{Joint normality.} Hence, up to $o_P(1)$, every $S_j(\phi)$ is a linear form
$\sum_{l\le K}c^{(n)}_{l}(E_l-1)$ in the same independent variables with coefficients bounded by
$C k^{-1/2}$. For any fixed linear combination of finitely many such forms, the coefficients are
again $O(k^{-1/2})$, so the Lindeberg condition holds and the Lindeberg--Feller theorem gives
asymptotic normality of the combination, provided its variance converges; the Cram\'er--Wold
device then gives joint asymptotic normality. Since the remainders and the deterministic terms
vanish in $L^2$, the variance of each combination differs by $o(1)$ from that of the same
combination of the $S_j(\phi)$, which converges by the covariance computation below.

\emph{Covariance.} Let $j<j'$. The excesses $\{V^{(j)}_i\}_{i\le k_j}$ are functions of
$E_1,\dots,E_{k_j}$, and $V^{(j')}_i=V^{(j)}_i+D$ for $i\le k_j$ with $D=\sum_{l=k_j+1}^{k_{j'}}E_l/l$,
while $V^{(j')}_i$ for $k_j<i\le k_{j'}$ depends only on $E_{k_j+1},\dots,E_{k_{j'}}$. Hence
$S_j(\phi)$ is independent of $D$ and of the second group, and
\[
\operatorname{cov}\{S_j(\phi),S_{j'}(\chi)\}=(k_jk_{j'})^{-1/2}\,\E_D\Bigl[\operatorname{cov}\Bigl\{\textstyle\sum_{i\le k_j}\phi(V^{(j)}_i),\sum_{i\le k_j}\chi(V^{(j)}_i+D)\Bigm|D\Bigr\}\Bigr].
\]
The multiset $\{V^{(j)}_i\}_{i\le k_j}$ has the law of an independent sample of $k_j$ standard
exponentials (R\'enyi), and both sums are symmetric functions of it, so the conditional
covariance equals $k_j\operatorname{cov}\{\phi(E),\chi(E+D)\}$ evaluated at the realised $D$. As
$D\to\log(k_{j'}/k_j)\to\log(t_{j'}/t_j)$ in probability and the map
$d\mapsto\operatorname{cov}\{\phi(E),\chi(E+d)\}$ is bounded and continuous (by the bounds on
$\chi'$), dominated convergence gives the stated limit. The case $j=j'$ gives
$\operatorname{var}S_j(\phi)=\operatorname{var}\phi(E)$ exactly. \qedhere
\end{proof}

\subsection{Local asymptotic normality}

Let $D^{(m)}=(X_{n-i+1,n}/X_{n-K,n})_{i\le K}$, taken in random order. Every statistic used in
this paper, at every grid point and working value, is a function of $D^{(m)}$. Given
$X_{n-K,n}=u$, $D^{(m)}$ is an independent sample from $q_u$; under exact Pareto tails it is an
independent sample from $p$.

\begin{lemma}\label{lem:lan}
Assume \ref{ass:A1}--\ref{ass:A3}. For any deterministic sequence $u_n$ with
$\sqrt{K}\,\delta(u_n)\to\lambda_m=\lambda$, the log-likelihood ratio of $q_{u_n}^{\otimes K}$
with respect to $p^{\otimes K}$ satisfies, under $p^{\otimes K}$,
$\log\prod_{i\le K}(q_{u_n}/p)(Y_i)=\lambda\Delta_K-\tfrac12\lambda^2\sigma_h^2+o_P(1)$, with
$\Delta_K=K^{-1/2}\sum_{i\le K}h_\rho(Y_i)\dto N(0,\sigma_h^2)$ and $\sigma_h^2=\int h_\rho^2p=\xi^{-2}\rho^2/(1-2\rho)$. The two
sequences of laws are mutually contiguous.
\end{lemma}

\begin{proof}
Put $\eta_n=\delta(u_n)$ and $W_{i}=2\{(q_{u_n}/p)^{1/2}(Y_i)-1\}$. The function $h_\rho$ satisfies $\int h_\rho p=0$ and $\int h_\rho^2p<\infty$. Both $q_{u_n}$ and $p$ are densities on $[1,\infty)$ and $p>0$ there, so $q_{u_n}$ is absolutely continuous with respect to $p$ and the likelihood ratio $q_{u_n}/p$ is well defined; we do not assume the converse, since $q_{u_n}$ may vanish on part of $[1,\infty)$. On $N_n=\{q_{u_n}=0\}$ the integrand of Assumption~\ref{ass:A2} equals $p(1+\frac12\eta_nh_\rho)^2\ge p/4$ for $n$ large, because $h_\rho$ is bounded; hence $\int_{N_n}p=o(\eta_n^2)$, and the probability under $p^{\otimes K}$ that some $Y_i$ falls in $N_n$ is at most $K\int_{N_n}p=o(K\eta_n^2)=o(1)$. So with probability tending to one every $W_i>-2$ and the log-likelihood ratio below is finite. The argument in the proof of Theorem~7.2 of \citet{vandervaart1998} uses only the expansion
$\int\{q_n^{1/2}-p^{1/2}-\frac12\eta_nh\,p^{1/2}\}^2=o(\eta_n^2)$ along the sequence and
$\sqrt{K}\eta_n\to\lambda$, and yields $\sum_iW_i=\lambda\Delta_K-\lambda^2\sigma_h^2/4+o_P(1)$,
$\sum_iW_i^2\to\lambda^2\sigma_h^2$ in probability and $\max_i|W_i|\to0$ in probability. Then
$\log\prod(q/p)=\sum_i2\log(1+W_i/2)=\sum_iW_i-\frac14\sum_iW_i^2+o_P(1)$. Asymptotic
normality of $\Delta_K$ is the central limit theorem, and contiguity follows from Le Cam's first
lemma because the limit of the log-likelihood ratio is $N(-\lambda^2\sigma_h^2/2,\lambda^2\sigma_h^2)$.
\qedhere
\end{proof}

Because $F$ is continuous above $x_0$, given $X_{n-K,n}=u>x_0$ the $K$ largest observations are
an independent sample from $F$ conditioned on exceeding $u$; this defines the regular conditional
law $\Prob_u$ of $D^{(m)}$. Since $X_{n-K,n}/U(n/K)\to1$ in probability and $\delta$ is regularly
varying, the uniform convergence theorem gives $\sqrt K\delta(u_n)\to\lambda$ for every sequence
with $u_n/U(n/K)\to1$.

\emph{Unconditioning argument.} Let $T_n$ be a function of $D^{(m)}$ and $B$ a set with
$L(\partial B)=0$, and suppose $\Prob_{u_n}(T_n\in B)\to L(B)$ for every deterministic sequence
$u_n$ with $u_n/U(n/K)\to1$. Put
$g_n(\eta)=\sup\{|\Prob_u(T_n\in B)-L(B)|:|u/U(n/K)-1|\le\eta\}$. Then $g_n(\eta_n)\to0$ for
every $\eta_n\to0$ (otherwise choosing near-maximising $u_n$ would contradict the hypothesis).
Choose $\eta_n\to0$ so slowly that $\Prob(|X_{n-K,n}/U(n/K)-1|>\eta_n)\to0$. Then
$|\Prob(T_n\in B)-L(B)|\le g_n(\eta_n)+2\Prob(|X_{n-K,n}/U(n/K)-1|>\eta_n)\to0$. The same
argument applies to expectations of bounded continuous functions, hence to convergence in
distribution and in probability. We use it below without further comment.


\begin{lemma}\label{lem:tests}
Fix $\rho'<0$, let $\Theta$ be a compact subset of the validity region of the EPD$(\rho')$ model
on which $1+\delta>0$, and let $\theta_0\in\Theta$. For a sample $Y_1,\dots,Y_k$ that is independent
with distribution EPD$(\theta;\rho')$, $\theta=(\xi,\delta)\in\Theta$, and for every $\epsilon>0$,
there are tests $\phi_k$ with $\E_{\theta_0}\phi_k\to0$ and
$\sup\{\E_\theta(1-\phi_k):\theta\in\Theta,\ \|\theta-\theta_0\|\ge\epsilon\}\to0$.
\end{lemma}

\begin{proof}
Write $G_\theta$ for the distribution function and $d(\theta,\theta')=\sup_{y\ge1}|G_\theta(y)-G_{\theta'}(y)|$.

\emph{Continuity of $d$.} Since $y^{\rho'/\xi}\in(0,1]$, the factor $1+\delta-\delta y^{\rho'/\xi}$
lies between $\min(1,1+\delta)$ and $\max(1,1+\delta)$, and by compactness
$\min_\Theta\min(1,1+\delta)=:b>0$ and $\max_\Theta\xi=:\bar\xi<\infty$. Hence
$1-G_\theta(y)\le(by)^{-1/\bar\xi}$ for $by\ge1$, uniformly over $\Theta$, so the tails beyond
a large $Y$ contribute uniformly little to $d$. On $\Theta\times[1,Y]$ the map
$(\theta,y)\mapsto G_\theta(y)$ is continuous, hence uniformly continuous. Together these give
$d(\theta_n,\theta)\to0$ whenever $\theta_n\to\theta$ in $\Theta$, so by the triangle inequality
$d$ is continuous on $\Theta\times\Theta$.

\emph{Identifiability.} As $y\to\infty$,
$\log\{1-G_\theta(y)\}=-\xi^{-1}\log y-\xi^{-1}\log(1+\delta)+o(1)$. If $G_\theta=G_{\theta'}$, the
coefficients of $\log y$ give $\xi=\xi'$, and the constants then give $\delta=\delta'$. So
$d(\theta,\theta')>0$ for $\theta\ne\theta'$.

\emph{Separation.} The set $\{(\theta,\theta')\in\Theta^2:\|\theta-\theta'\|\ge\epsilon/2\}$ is
compact and $d$ is continuous and positive on it, so
$\eta=\min\{d(\theta,\theta'):\|\theta-\theta'\|\ge\epsilon/2\}>0$ (if the set is empty there is
nothing to prove).

\emph{Minimum-distance estimator and test.} Let $F_k$ be the empirical distribution function and
$d_k(\theta)=\sup_y|F_k(y)-G_\theta(y)|$. Since $|d_k(\theta)-d_k(\theta')|\le d(\theta,\theta')$,
$d_k$ is continuous on the compact $\Theta$ and has a minimiser $\hat\theta_k$. By the
Dvoretzky--Kiefer--Wolfowitz inequality with the constant of \citet{massart1990},
$\Prob_\theta\{d_k(\theta)>\eta/3\}\le2e^{-2k\eta^2/9}$ for every $\theta$. On the complement,
$d(\hat\theta_k,\theta)\le d_k(\hat\theta_k)+d_k(\theta)\le2d_k(\theta)\le2\eta/3<\eta$, hence
$\|\hat\theta_k-\theta\|<\epsilon/2$. Put $\phi_k=1\{\|\hat\theta_k-\theta_0\|\ge\epsilon/2\}$. Then
$\E_{\theta_0}\phi_k\le2e^{-2k\eta^2/9}$, and if $\|\theta-\theta_0\|\ge\epsilon$, then
$\|\hat\theta_k-\theta\|<\epsilon/2$ forces $\|\hat\theta_k-\theta_0\|>\epsilon/2$, so
$\E_\theta(1-\phi_k)\le2e^{-2k\eta^2/9}$ uniformly in such $\theta$. \qedhere
\end{proof}

\subsection{Proofs of the main results}

\begin{proof}[Proof of Theorem~\ref{thm:bvm}]
Fix a sequence $u_n$ as in Lemma~\ref{lem:lan}; all probability statements are first conditional
on $X_{n-K,n}=u_n$ and then unconditioned.

(ii) Under $p^{\otimes K}$, by the remark after Lemma~\ref{lem:clt},
$T_{jl}:=\sqrt{k_j}\{\tilde\xi_j(\rho_l)-\xi\}/\xi=S_j(\psi_{\rho_l})$ and
$\xi\Delta_K=S_m(a_\rho)$, jointly. Hence $(T_{jl},\Delta_K)$ is jointly asymptotically normal under
$p^{\otimes K}$, and $T\dto G$, which defines $G$ and gives \eqref{eq:cov}. By Lemma~\ref{lem:clt}
with $j'=m$ and $\chi=a_\rho$, $\lim\operatorname{cov}(T_{jl},\xi\Delta_K)=t_j^{1/2}
\operatorname{cov}\{\psi_{\rho_l}(E),a_\rho(E+\log t_j^{-1})\}$. Since
$a_\rho(E+r)=e^{\rho r}a_\rho(E)+(e^{\rho r}-1)$ and $\psi_{\rho_l}$ has mean zero, this equals
$t_j^{1/2-\rho}\operatorname{cov}\{\psi_{\rho_l}(E),a_\rho(E)\}=t_j^{1/2-\rho}\beta(\rho,\rho_l)$
by Lemma~\ref{lem:scores}(iv). Le Cam's third lemma \citep[Example~6.7]{vandervaart1998} and
Lemma~\ref{lem:lan} give, under $q_{u_n}^{\otimes K}$, $T\dto G+\mu$ with
$\mu_{jl}=\lambda t_j^{1/2-\rho}\beta(\rho,\rho_l)/\xi=(\lambda_j/\xi)\beta(\rho,\rho_l)$.

(iii) Under $p^{\otimes K}$, Lemma~\ref{lem:scores}(v) applied at threshold $k_j$ gives
$\sqrt{k_j}(\hat\xi_j-\tilde\xi_j)\to0$ in probability; contiguity transfers this to
$q_{u_n}^{\otimes K}$.

(i) Under $p^{\otimes K}$, the excesses over $X_{n-k_j,n}$ are, given that order statistic, an
independent sample of size $k_j$ from the EPD$(\xi,0;\rho_l)$ model. This model is
differentiable in quadratic mean at $(\xi,0)$, because on a neighbourhood of $(\xi,0)$ inside the
open validity region $\{\delta>\max(-1,\xi/\rho_l)\}$ its density is continuously
differentiable in the parameters with continuous Fisher information
\citep[Lemma~7.6]{vandervaart1998}; the information is nonsingular by
Lemma~\ref{lem:scores}(i). Uniformly consistent tests of $(\xi,0)$ against $\{\|\theta-(\xi,0)\|\ge\epsilon\}$ within the
support of the prior exist by Lemma~\ref{lem:tests} below. The Bernstein--von Mises theorem
\citep[Theorem~10.1]{vandervaart1998} gives total-variation convergence of the posterior of
$\sqrt{k_j}(\vartheta-\xi,\delta)$ to the normal law with mean the efficient statistic and
covariance the inverse information; its $\xi$-marginal is $N\{\sqrt{k_j}(\tilde\xi_j-\xi),
\xi^2s(\rho_l)^2\}$ by Lemma~\ref{lem:scores}(ii). The total-variation distance is a function of
$D^{(m)}$ that tends to zero in probability under $p^{\otimes K}$, hence under
$q_{u_n}^{\otimes K}$ by contiguity. \qedhere
\end{proof}

\begin{proof}[Proof of Corollary~\ref{cor:known}]
By Theorem~\ref{thm:bvm}(i) and continuity of the normal distribution function, the endpoints of
$B_j(\rho)$ are $\tilde\xi_j(\rho)\pm c(\rho)\xi s(\rho)/\sqrt{k_j}+o_P(k_j^{-1/2})$. The
exact $(1-\alpha)$-quantile of the pivot under Pareto tails at the actual grid converges to
$c(\rho)$, because by Lemmas~\ref{lem:clt} and~\ref{lem:scores}(v) the pivot converges in law to $\max_j|G_j|/s(\rho)$, whose
distribution function is continuous and strictly increasing on $(0,\infty)$ for a non-degenerate
Gaussian vector. (The computed constant is a Monte Carlo estimate of this quantile; the theory
refers to the exact quantile, or to a Monte Carlo size that grows with $n$.) Hence $\{\xi\in B_j(\rho)\ \forall j\}$ differs from
$\{\max_j|T_{j}|\le c(\rho)s(\rho)+o_P(1)\}$ by an event of vanishing probability, and with
$\mu=0$ (since $\beta(\rho,\rho)=0$) the probability tends to
$\Prob\{\max_j|G_j|\le c(\rho)s(\rho)\}=1-\alpha$, the distribution of $\max_j|G_j|$ being
continuous. \qedhere
\end{proof}

\begin{proof}[Proof of Proposition~\ref{prop:firstorder}]
The Pareto posterior is the posterior in the one-parameter Pareto model, with a prior on $\xi$
that is continuous and positive at $\xi$ and compactly supported (the one-parameter analogue of
Assumption~\ref{ass:P}); the argument of Theorem~\ref{thm:bvm} with score $E-1$ shows that its centre is the Hill estimator with
$\sqrt{k_j}(H_j-\xi)/\xi\dto G^H_j+\mu^H_j$, $\mu^H_j=(\lambda_j/\xi)\rho/(1-\rho)$, where $G^H$ has
covariance $(t_j/t_{j'})^{1/2}$ by Lemma~\ref{lem:clt}, which is nonsingular for distinct $t_j$.
The limiting coverage is $f(1)$ with $f(t)=\Prob(G^H+t\mu^H\in C)$ and
$C=\{g:\max_j|g_j|\le c^H\}$, a bounded symmetric convex set with nonempty interior. By
\citet{anderson1955}, $f$ is non-increasing on $[0,\infty)$. It is real-analytic, being an
integral over $C$ of the non-degenerate Gaussian density shifted by $t\mu^H$, and $f(t)\to0$ as $t\to\infty$
because $C$ is bounded. If $f(1)=f(0)$, $f$ would be constant on $[0,1]$ and hence everywhere,
which is impossible. So $f(1)<f(0)=1-\alpha$. As $|\lambda|\to\infty$, $|\mu^H_m|\to\infty$ and
$f(1)\le\Prob(|G^H_m+\mu^H_m|\le c^H)\to0$. \qedhere
\end{proof}

\begin{proof}[Proof of Theorem~\ref{thm:union}]
(a) As in the proof of Corollary~\ref{cor:known}, for each $l$ the endpoints of $B_j(\rho_l)$ are
$\tilde\xi_j(\rho_l)\pm c_l\xi s_l/\sqrt{k_j}+o_P(k_j^{-1/2})$, so after multiplying by
$\sqrt{k_j}/\xi$ the coverage event becomes
$\{\forall j:\min_l(T_{jl}-c_ls_l)\le o_P(1),\ \max_l(T_{jl}+c_ls_l)\ge o_P(1)\}$. By
Theorem~\ref{thm:bvm}(ii) and the continuous-mapping theorem, the probability converges to
$\Gamma_\Mesh(\rho,\lambda)$, because the boundary of the limiting event is contained in the
finite union of sets $\{g_{jl}=\pm c_ls_l-\mu_{jl}\}$, each of probability zero since every $G_{jl}$
has variance $s_l^2>0$.
(b) If $\rho=\rho_l$ then $\mu_{jl}=0$ for all $j$, and the limiting event contains
$\{\max_j|G_{jl}|\le c_ls_l\}$, of probability $1-\alpha$. If $\lambda=0$ then $\mu=0$ and the same
holds for every $l$.
(c) $\{\xi\in U_j\ \forall j\}\subset\{\xi\in U_{\hat\jmath}\}$ pointwise on the sample space.
(d) By Lemma~\ref{lem:beta}, $\beta(\rho,\rho_l)<0<\beta(\rho,\rho_{l'})$. Let $\lambda\to+\infty$
(the case $-\infty$ is symmetric). Every $\lambda_j=\lambda t_j^{1/2-\rho}\to+\infty$, so
$\mu_{jl}\to-\infty$ and $\mu_{jl'}\to+\infty$ for every $j$, and
$\Gamma_\Mesh\ge1-\sum_j\{\Prob(G_{jl}+\mu_{jl}-c_ls_l>0)+\Prob(G_{jl'}+\mu_{jl'}+c_{l'}s_{l'}<0)\}\to1$.
(e) If $\rho<\rho_1$, all $\beta(\rho,\rho_l)>0$. As $\lambda\to+\infty$,
$\min_l(G_{ml}+\mu_{ml}-c_ls_l)\to+\infty$ in probability, so the limiting event at $j=m$, and
hence $\Gamma_\Mesh$, has probability tending to zero; $\lambda\to-\infty$ and $\rho>\rho_L$ are
symmetric.
(f) The first statement follows as in (a). The hull contains every $B_j(\rho_l)$, whose scaled
width tends to $2c_ls_l$, which gives the lower bound. Dividing by $|\lambda_j|$, $G/|\lambda_j|\to0$
and $\mu_{jl}/|\lambda_j|=\pm\beta(\rho,\rho_l)/\xi$ give the last statement. \qedhere
\end{proof}

\begin{proof}[Proof of Theorem~\ref{thm:uniform}]
Write $\Gamma_\Mesh=1-\Prob(\mathcal{F})$, where $\mathcal{F}=\bigcup_j(F^+_j\cup F^-_j)$ with
$F^+_j=\{\min_l(G_{jl}+\mu_{jl}-c_ls_l)>0\}$ (the hull lies above $\xi$ at $k_j$) and
$F^-_j=\{\max_l(G_{jl}+\mu_{jl}+c_ls_l)<0\}$ (it lies below). Taking $l=1$ in the minimum and
$l=L$ in the maximum,
\[
F^+_j\subseteq\{Z_{j1}>c_1-\mu_{j1}/s_1\},\qquad F^-_j\subseteq\{Z_{jL}<-c_L-\mu_{jL}/s_L\}.
\]
Let $\rho\in R$ and $\lambda\ge0$. Then $\lambda_j=\lambda t_j^{1/2-\rho}\ge0$, and since
$\rho_1\le\rho\le\rho_L$, Lemma~\ref{lem:beta} gives $\beta(\rho,\rho_1)\le0\le\beta(\rho,\rho_L)$,
so $\mu_{j1}\le0\le\mu_{jL}$ for every $j$. Hence $F^+_j\subseteq\{Z_{j1}>c_1\}$ and
$F^-_j\subseteq\{Z_{jL}<-c_L\}$, and $\Prob(\mathcal{F})\le\varphi_\Mesh$. If $\lambda<0$ the signs are
reversed, and the same two inclusions with $l=L$ in the minimum and $l=1$ in the maximum give
$\mathcal{F}\subseteq\bigcup_j\{Z_{jL}>c_L\}\cup\{Z_{j1}<-c_1\}$, whose probability equals $\varphi_\Mesh$
because $G$ and $-G$ have the same law. The union bound and the same symmetry give
$\varphi_\Mesh\le\pi_1+\Prob(\min_jZ_{jL}<-c_L)=\pi_1+\pi_L$. The statement about $\hat\jmath$
follows from Theorem~\ref{thm:union}(a) and (c). \qedhere
\end{proof}

\begin{proof}[Proof of Corollary~\ref{cor:uniform}]
For the two-sided constants, $\alpha=\Prob(\max_j|Z_{jl}|>c_l)=\Prob(\max_jZ_{jl}>c_l)+
\Prob(\min_jZ_{jl}<-c_l)-q_l=2\pi_l-q_l$, using the symmetry of $G$; so $\pi_l=(\alpha+q_l)/2$ and
(ii) follows from Theorem~\ref{thm:uniform}. For (i), $\Prob(\max_jZ_{jl}>c_l)=\pi_l\ge\alpha/2$, so
$c_l^+\ge c_l$ and $U_j\subseteq U_j^+$. The proof of Theorem~\ref{thm:union}(a) applies to $U^+$
with $c_1,c_L$ replaced by $c_1^+,c_L^+$ (the distribution functions of $\max_jZ_{jl}$ and $\max_j(-Z_{jl})$ are equal, continuous
and strictly increasing, so both finite-grid one-sided quantiles, and hence their maximum, converge
to $c_l^+$ as in the proof of Corollary~\ref{cor:known}), and the proof of Theorem~\ref{thm:uniform} with these constants gives
a limiting non-coverage probability at most $\Prob(\max_jZ_{j1}>c_1^+)+\Prob(\min_jZ_{jL}<-c_L^+)=\alpha$.
Finally $\Prob(\max_j|Z_{jl}|>c_l^+)\le2\Prob(\max_jZ_{jl}>c^+_l)=\alpha$, so each endpoint band with
constant $c_l^+$ is a valid two-sided simultaneous band. \qedhere
\end{proof}

\begin{proof}[Proof of Proposition~\ref{prop:tight}]
Let $\rho=\rho_1$ and $\lambda\to+\infty$. Then $\mu_{j1}=0$, while for $l\ge2$,
$\beta(\rho_1,\rho_l)>0$ by Lemma~\ref{lem:beta} and $\mu_{jl}\to+\infty$ for every $j$. Hence,
almost surely, $\min_l(G_{jl}+\mu_{jl}-c_ls_l)\to G_{j1}-c_1s_1$ and
$\max_l(G_{jl}+\mu_{jl}+c_ls_l)\to+\infty$, so the indicator of the coverage event converges to
the indicator of $\{\max_jZ_{j1}\le c_1\}$ except on the null set $\{\max_jZ_{j1}=c_1\}$. By
dominated convergence $\Gamma_\Mesh(\rho_1,\lambda)\to1-\pi_1$. For $\rho=\rho_L$ and
$\lambda\to-\infty$, $\mu_{jL}=0$ and, for $l<L$, $\beta(\rho_L,\rho_l)<0$, so again $\mu_{jl}\to+\infty$;
the same argument gives $\Gamma_\Mesh(\rho_L,\lambda)\to1-\Prob(\max_jZ_{jL}>c_L)=1-\pi_L$. Since
$\rho_1,\rho_L\in R$, the infimum over $R\times\R$ is at most both limits; the lower bound is
Theorem~\ref{thm:uniform}. \qedhere
\end{proof}

\begin{proof}[Proof of Proposition~\ref{prop:cex}]
Let $X,Y$ be independent standard normal, $(Z_{j1})_{j\le4}=(X,-X,X,-X)$ and
$(Z_{j2})_{j\le4}=(Y,-Y,Y,-Y)$. Then $\max_j|Z_{jl}|$ is $|X|$ or $|Y|$, so the two-sided constants
are $c_1=c_2=c:=z_{1-\alpha/2}$. Take $\nu_{\cdot1}=(0,0,-b,-b)$ and $\nu_{\cdot2}=(b,b,0,0)$ with
$b>0$, which satisfy $\nu_{j1}\le0\le\nu_{j2}$. The hull fails at $j$ if both shifted components
exceed $c$ or both are below $-c$. As $b\to\infty$ these events converge almost surely to
$\{X>c\}$ ($j=1$), $\{X<-c\}$ ($j=2$), $\{Y<-c\}$ ($j=3$) and $\{Y>c\}$ ($j=4$), because the
component carrying the shift $\pm b$ is eventually above $c$ or below $-c$, respectively. So the
probability of the hull event tends to $\Prob(|X|\le c,|Y|\le c)=(1-\alpha)^2$. Here
$q_l=\Prob(\max_jZ_{jl}>c,\min_jZ_{jl}<-c)=\Prob(|X|>c)=\alpha$. \qedhere
\end{proof}

\begin{proof}[Proof of Corollary~\ref{cor:intersection}]
$\xi\in\bigcap_jU_j$ if and only if $\xi\in U_j$ for every $j$. The width bound is trivial, and
$\{I=\emptyset\}\subset\{\xi\notin I\}$. \qedhere
\end{proof}

\begin{proof}[Proof of Proposition~\ref{prop:exact}]
On $\{X_{n-K,n}>x_0\}$ the relative excesses above $X_{n-K,n}$ are an independent Pareto sample
(as a set), so by R\'enyi's representation the vector $(\hat\xi_j(\rho')/\xi)_{j\le m}$ has a law
that depends only on the grid and $\rho'$, and so does
$Q=\max_j\sqrt{k_j}|\hat\xi_j(\rho')/\xi-1|/s(\rho')$. $Q$ has a continuous distribution: each $\hat\xi_j(\rho')$ is a real-analytic function of $(E_l)$ on the open positive orthant and is positively homogeneous of degree one, hence non-constant, so its level sets have Lebesgue measure zero, and $\Prob(Q=c)\le\sum_j\Prob\{\sqrt{k_j}|\hat\xi_j(\rho')/\xi-1|=c\,s(\rho')\}=0$. If $c$ is the $(1-\alpha)$-quantile of $Q$,
then $\Prob(Q\le c)=1-\alpha$, and $\{Q\le c\}$ is the event that $\xi$ lies in every interval of
the stated band. \qedhere
\end{proof}

\subsection{The covariance of the limit}\label{app:cov}

Write $\psi_r(E)=A_r(E-1)-B_ra_r(E)$ with $A_r=(1-r)^2/r^2$ and $B_r=(1-2r)(1-r)/r^3$. For $d\ge0$,
$a_r(E+d)=(1-r)e^{rd}e^{rE}-1$, and with $E\sim\mathrm{Exp}(1)$, $\operatorname{var}E=1$,
$\operatorname{cov}(E,e^{sE})=s/(1-s)^2$ and
$\operatorname{cov}(e^{sE},e^{uE})=1/(1-s-u)-1/\{(1-s)(1-u)\}$ for $s,u<1$ with $s+u<1$ (here $s,u\le0$), so
\begin{multline*}
\operatorname{cov}\{\psi_r(E),\psi_{r'}(E+d)\}=A_rA_{r'}-A_rB_{r'}e^{r'd}\frac{r'}{1-r'}
-A_{r'}B_r\frac{r}{1-r}\\
+B_rB_{r'}(1-r)(1-r')e^{r'd}\Bigl\{\frac1{1-r-r'}-\frac1{(1-r)(1-r')}\Bigr\}.
\end{multline*}
With \eqref{eq:cov} this gives the covariance matrix of $G$ in closed form; at $d=0$, $r=r'$ it
reduces to $s(r)^2$, as Lemma~\ref{lem:scores}(ii) requires. The numerical evaluations of
$\Gamma_\Mesh$, $\pi_l$, $q_l$ and $\varphi_\Mesh$ in Remark~\ref{rem:numeric} and
Figure~\ref{fig:limit} sample $G$ from this Gaussian law; the prediction check of
Section~\ref{sec:predict} instead simulates the estimators under Pareto tails at each setting's own
grid, which retains finite-$k$ effects.

\section{Reproducibility and data}\label{app:repro}

All results are generated by scripts in the repository; seeds are recorded in the scripts and
reports. The pre-registered benchmark (code frozen before the run, raw outputs and file hashes
stored) uses a single test seed. Post hoc analyses use development seeds or regenerate the
same benchmark samples from the stored seed, and are labelled as post hoc. The Danish fire
losses are read at run time from the CRAN mirror of the R package \texttt{evir} (GPL $\ge2$) and
are not redistributed. No other data are used.

\begin{table}[ht]
\centering\small
\caption{Analyses, their status, and their sources in the repository.}
\label{tab:analyses}
\begin{tabular}{lll}
\toprule
Analysis & Status & Script / output\\
\midrule
Benchmark, Table~\ref{tab:bench} & pre-registered & \texttt{scripts/run\_phase3\_benchmark.py}\\
Limit $\Gamma_\Mesh$, Figure~\ref{fig:limit} & post hoc (theory) & \texttt{scripts/theory\_union\_limit.py}\\
$\pi_l$, $q_l$, $\varphi_\Mesh$, $U^+$, Remark~\ref{rem:numeric} & post hoc (theory) & \texttt{scripts/theory\_union\_endpoint\_bound.py}\\
Prediction check, Figure~\ref{fig:predict} & post hoc & \texttt{scripts/theory\_union\_predict\_benchmark.py}\\
Critical-value rescaling & post hoc & \texttt{scripts/theory\_union\_kappa.py}\\
$\hat\rho$ audit & post hoc & \texttt{scripts/dev\_phase3\_rhohat\_tuning.py}\\
Comparators, Table~\ref{tab:posthoc} & post hoc & \texttt{scripts/posthoc\_comparators.py}\\
Danish data, Table~\ref{tab:danish} & frozen method & \texttt{scripts/apply\_danish.py}\\
Numbers, tables, figures & --- & \texttt{scripts/build\_preprint\_assets.py}\\
\bottomrule
\end{tabular}
\end{table}
